\documentclass[11pt]{article}

\usepackage[a4paper,margin=1in]{geometry}
\usepackage{amsmath, amssymb}
\usepackage{microtype}
\usepackage{newpxtext, newpxmath}
\usepackage{xcolor}
\usepackage[most]{tcolorbox}
\usepackage{enumitem}
\usepackage{titlesec}

\definecolor{inkblue}{HTML}{1F3A5F}
\definecolor{boxbg}{HTML}{F2F5F9}
\definecolor{boxrule}{HTML}{8AA4C4}

\titleformat{\section}{\large\bfseries\color{inkblue}}{}{0em}{}
\titlespacing{\section}{0pt}{1.4em}{0.5em}

\newtcolorbox{insight}{
  colback=boxbg, colframe=boxrule, boxrule=0.6pt,
  arc=2pt, left=10pt, right=10pt, top=8pt, bottom=8pt,
  breakable, enhanced,
  title={\textbf{Insight}}, coltitle=inkblue, fonttitle=\bfseries\small,
  attach title to upper=\hspace{0.6em}
}

\setlength{\parskip}{0.55em}
\setlength{\parindent}{0pt}

\title{\color{inkblue}\bfseries Next-Token Prediction \textit{is} Geometry}
\author{A bridge between conditional probability and high-dimensional manifolds}
\date{}

\begin{document}
\maketitle
\thispagestyle{empty}

One could sum up this whole document in a sentence:
\begin{tcolorbox}[colback=yellow!18, colframe=yellow!55!black, arc=2pt, boxrule=0.5pt,
  left=10pt, right=10pt, top=8pt, bottom=8pt]
A large language model learns a map whose image is a sub-manifold of a very
high-dimensional ambient space, on which geometry tracks meaning.
\end{tcolorbox}
Arbitrary gibberish gets mapped to arbitrary points in the ambient space, but
\emph{meaningful} language gets mapped to a lower-dimensional sub-manifold.

We are using the term \emph{manifold} loosely: the set of
vectors the model produces need not be smooth, need not have constant dimension, and may be
singular or stratified.

\section{Next token prediction}

A large language model is a function that takes a context and returns a probability
distribution over the next token:
\[
  P(\text{next} \mid \text{context}).
\]
That is the whole type signature, and the ``statistical next-token prediction'' view
has it exactly right. Sampling from that distribution, one token at a time, is the
entire performance.

\section{The one equation that is the bridge}

Inside the model, the context is compressed to a single vector $h \in \mathbb{R}^d$,
with $d$ a few thousand. Every candidate token $w$ carries its own vector, its embedding
$e_w$, living in the \emph{same} space. The probability we want is then
\[
  P(w \mid \text{context})
  \;=\; \operatorname{softmax}_w \langle h,\, e_w \rangle
  \;=\; \frac{\exp\langle h,\, e_w\rangle}{\displaystyle\sum_{w'} \exp\langle h,\, e_{w'}\rangle}.
\]
The inner product $\langle h, e_w\rangle$ is nothing but \emph{how aligned the context's
direction is with the token's direction} --- an angle and a length. The softmax is
merely the rule that turns a list of alignments into a list of probabilities summing to
one. So:

\begin{quote}
  \textbf{The conditional probability is high exactly when the context vector points the
  same way as the token vector.}
\end{quote}

``Likely next word'' and ``nearby direction in high-dimensional space'' are not two
facts. They are one fact, written in two different notations.

\bigskip
\begin{center}\color{boxrule}\rule{0.42\linewidth}{0.4pt}\end{center}
\bigskip

\subsection*{Mathematics}

A language model works with a
fixed, finite list of symbols chosen before training --- the \emph{vocabulary} --- whose
$n$ entries are its \emph{tokens}. Each token $w$ carries its embedding
vector $e_w$, and predicting the next token means choosing one entry from this list.

The context --- the tokens so far --- is read by the model and compressed into a single
vector $h$ of the same dimension, so it lives in the same space as the token embeddings. It
has to: the model predicts by comparing $h$ against each token's vector and asking which it
best matches, and such a comparison only means something when the two vectors share a space.
So $h$ is built precisely to be measured against the $e_w$ --- the context rendered as a
point among the tokens.

Recall how the inner product of two vectors encodes the angle between them:
\[
  \langle u, v\rangle = \lVert u\rVert\,\lVert v\rVert\cos\theta,
\]
where $\theta$ is the angle separating them.

A \emph{logit} is the raw, un-normalised score the model gives a token before it becomes a
probability:
\[
  z_w = \langle h, e_w\rangle .
\]
The vector of all $n$ logits is what softmax turns into probabilities.

Recall two facts about the exponential function: it is always positive, $e^x > 0$, and it
is equal to its own derivative, $\tfrac{d}{dx}e^x = e^x$.

Given a vector
of logits $z = (z_1, \dots, z_n) \in \mathbb{R}^n$ --- one entry per vocabulary token --- the
\emph{softmax} is
\[
  \operatorname{softmax}(z)_i \;=\; \frac{e^{z_i}}{\sum_{k=1}^{n} e^{z_k}} \;=:\; p_i .
\]

\paragraph{It is a probability distribution.} Two checks. Each numerator $e^{z_i}$ is
strictly positive and the denominator is a sum of positives, so $p_i > 0$ for every $i$;
and they normalise,
\[
  \sum_{i=1}^{n} p_i = \frac{\sum_{i=1}^n e^{z_i}}{\sum_{k=1}^n e^{z_k}} = 1 .
\]

Recall two facts about the logarithm: it is the inverse of the exponential, and it turns
products into sums, $\log(ab) = \log a + \log b$. As a consequence it turns quotients into
differences, $\log(a/b) = \log a - \log b$.

Now, writing $S = \sum_k e^{z_k}$,
\[
  \log p_w = z_w - \log S .
\]
Differences of logits shed that constant and give us \emph{log-odds},
\[
  z_i - z_j = \log\frac{p_i}{p_j} .
\]

Recall the quotient rule from calculus: for differentiable functions $f$ and $g$,
\[
  \left(\frac{f}{g}\right)' = \frac{f'g - f g'}{g^2}.
\]

On the diagonal ($i = j$), the quotient rule gives
\[
  \frac{\partial p_i}{\partial z_i}
    = \frac{e^{z_i}\,S - e^{z_i}\,e^{z_i}}{S^2}
    = p_i - p_i^2 = p_i\,(1 - p_i).
\]
Off the diagonal ($i \neq j$),
\[
  \frac{\partial p_i}{\partial z_j}
    = e^{z_i}\cdot\Bigl(-\frac{1}{S^2}\Bigr)e^{z_j}
    = -\,p_i\,p_j .
\]

\subsection*{The loss}

Recall two additional facts about the logarithm. $\log(1) = 0$ and $\log(x)$ tends to
negative infinity as $x$ tends to $0$.

Define the log loss by
\[
  L = -\log p_y .
\]
If the model is certain that the next token is $y$ then the loss is zero. If the model has
a low probability of predicting the correct next token then the loss is large.

The corpus the model trains on is a collection of (context, next-token) pairs. The total loss is the sum of the per-pair
losses,
\[
  L_{\text{total}} = \sum_t -\log p_{y_t} = -\log \prod_t p_{y_t},
\]
The product $\prod_t p_{y_t}$ is the probability of independent events. Minimizing the total
loss is the same as maximizing $\prod_t p_{y_t}$. Maximizing the probability that you will
guess the correct next token over the entire corpus.

\subsection*{Differentiating the loss}

Recall the derivative of the logarithm:
\[
  \frac{d}{dx}\log x = \frac{1}{x}.
\]

Now, by the chain rule,
\[
  \frac{\partial L}{\partial z_j}
    = \frac{\partial L}{\partial p_y}\,\frac{\partial p_y}{\partial z_j}
    = -\frac{1}{p_y}\,\frac{\partial p_y}{\partial z_j}.
\]
Now use the derivatives of the softmax that we calculated earlier. For $j = y$,
\[
  \frac{\partial L}{\partial z_j} = -\frac{1}{p_y}\,p_y(1-p_y) = -(1-p_y) = p_y - 1 = p_j - 1 .
\]
For $j \neq y$,
\[
  \frac{\partial L}{\partial z_j} = -\frac{1}{p_y}\,(-p_y p_j) = p_j .
\]
We combine these two facts using the Kronecker delta notation:
\[
  \frac{\partial L}{\partial z_j} = p_j - \delta_{jy}.
\]

\subsection*{The gradient is a spring law}

Recall that the logits are defined by the inner product $z_w = \langle h, e_w\rangle$.

The embedding of the context $h$ is a vector, so we expand it in the standard basis,
\[
  h = \sum_i h_i\,\hat{e}_i.
\] Written in coordinates, each logit is a dot product: expanding $h$ in the basis,
\[
  z_w = \langle h, e_w\rangle
      = \Bigl\langle \sum_i h_i\,\hat{e}_i,\; e_w \Bigr\rangle
      = \sum_i h_i\,\langle \hat{e}_i, e_w\rangle.
\]
So $z_w$ is a linear function of the $h_i$.

The loss is a function of these logits, so the chain rule gives
\[
  \frac{\partial L}{\partial h_i}
    = \sum_w \frac{\partial L}{\partial z_w}\,\frac{\partial z_w}{\partial h_i}
    = \sum_w \frac{\partial L}{\partial z_w}\, \langle \hat{e}_i, e_w\rangle.
\]
Stacking the coordinates back into a vector,
\[
  \frac{\partial L}{\partial h} = \sum_i \frac{\partial L}{\partial h_i}\,\hat{e}_i ,
\]
and swapping the order of the two sums,
\[
  \frac{\partial L}{\partial h}
    = \sum_i \sum_w \frac{\partial L}{\partial z_w}\,\langle \hat{e}_i, e_w\rangle\,\hat{e}_i
    = \sum_w \frac{\partial L}{\partial z_w} \Bigl( \sum_i \langle \hat{e}_i, e_w\rangle\,\hat{e}_i \Bigr)
    = \sum_w \frac{\partial L}{\partial z_w}\, e_w .
\]
So, with $\partial L/\partial z_w = p_w - \delta_{wy}$,
\[
  \frac{\partial L}{\partial h} = \sum_w (p_w - \delta_{wy})\, e_w = \bar{e} - e_y,
  \qquad \bar{e} = \sum_w p_w\, e_w .
\]

Similarly, expand
\[
  e_\omega = \sum_i \langle e_\omega, \hat{e}_i\rangle\,\hat{e}_i.
\]
so
\[
  \frac{\partial L}{\partial e_\omega}
    = \sum_i \frac{\partial L}{\partial \langle e_\omega, \hat{e}_i\rangle}\,\hat{e}_i .
\]
Recall that the loss function is a function of the logits $z_k$. By the chain rule:
\[
  \frac{\partial L}{\partial \langle e_\omega, \hat{e}_i\rangle}
    = \sum_k \frac{\partial L}{\partial z_k}\,\frac{\partial z_k}{\partial \langle e_\omega, \hat{e}_i\rangle}.
\]
Since $z_k$ depends only on $e_k$, if $k \neq \omega$ then varying the $i$-th component of
$e_\omega$ does not change $z_k$, so
\[
  \frac{\partial z_k}{\partial \langle e_\omega, \hat{e}_i\rangle} = 0 ;
\]
only the $k = \omega$ term survives:
\[
  \frac{\partial L}{\partial \langle e_\omega, \hat{e}_i\rangle}
    = \frac{\partial L}{\partial z_\omega}\,\frac{\partial z_\omega}{\partial \langle e_\omega, \hat{e}_i\rangle}
    = \frac{\partial L}{\partial z_\omega}\, h_i .
\]
Substituting,
\[
  \frac{\partial L}{\partial e_\omega}
    = \sum_i \frac{\partial L}{\partial z_\omega}\, h_i\,\hat{e}_i
    = \frac{\partial L}{\partial z_\omega}\sum_i h_i\,\hat{e}_i
    = \frac{\partial L}{\partial z_\omega}\, h
    = (p_\omega - \delta_{\omega y})\, h .
\]

\begin{tcolorbox}[colback=yellow!18, colframe=yellow!55!black, arc=2pt, boxrule=0.5pt,
  breakable, left=10pt, right=10pt, top=8pt, bottom=8pt]
Take one training step. A context produces the hidden vector $h$, the true next token is
$y$, and the model's current distribution is $p_w = \operatorname{softmax}\langle h, e_w\rangle$.
The loss is $L = -\log p_y$. Differentiating is the whole mechanism in two lines:
\[
  \frac{\partial L}{\partial h} = \sum_{w}\bigl(p_w - \delta_{wy}\bigr)e_w
    = \bar e - e_y,
  \qquad
  \frac{\partial L}{\partial e_w} = \bigl(p_w - \delta_{wy}\bigr)\,h,
  \qquad \bar e = \sum_w p_w\, e_w.
\]
Training moves against the gradient, so the update is its negative. It says four things ---
two about the context, two about the embeddings.
\begin{itemize}[leftmargin=1.4em]
  \item \textbf{The context moves toward the true token.} The update to $h$ is
  $-\partial L/\partial h = e_y - \bar e$; its $+e_y$ part pulls $h$ toward the embedding of
  the token that actually came next.
  \item \textbf{The context moves away from what it expected.} The same update's $-\bar e$
  part pushes $h$ away from $\bar e$, the average embedding under its current distribution.
  \item \textbf{The true token's embedding moves toward the context.} The update to $e_y$
  is $(1 - p_y)\,h$, pulling $e_y$ toward $h$.
  \item \textbf{Every other token embedding moves away from the context.} The update to $e_w$
  ($w \neq y$) is $-p_w\,h$, pushing it away from $h$ by exactly $p_w$ --- the probability
  wrongly placed on it.
\end{itemize}
\end{tcolorbox}

That is a force-directed layout. Every occurrence in the corpus is one tug: this context
attracts this token's direction and repels the others. Summed over billions of examples,
the attractions and repulsions reach an equilibrium --- and \textbf{the equilibrium
configuration is the manifold}.

\subsection*{Why the equilibrium is a \emph{meaning} manifold}

The output is a smooth, fixed function of $h$ alone. So the only way the model can give
two contexts similar next-token distributions is to map them to nearby $h$. Contexts
followed by similar things are forced together; contexts followed by different things are
forced apart. That is the distributional hypothesis --- \emph{you shall know a word by the
company it keeps} --- made mechanical.

\end{document}
